\documentclass[12pt]{article}
\usepackage{amsmath, amssymb}
\usepackage[margin=1in]{geometry}
\setlength{\parskip}{6pt}
\title{QUBO Formulation: Quantum Circuit Compilation Candidate Selection Problem}
\date{}
\begin{document}
\maketitle

\subsection*{Quantum Circuit Compilation Candidate Selection Problem}

\paragraph{Problem Statement.}
Given a fixed set of $N$ candidate compiled circuits, each associated with a precomputed scalar compilation cost $C_i \geq 0$ that aggregates depth, fidelity loss, and hardware constraint penalties, the task is to select exactly one candidate such that the total compilation cost is minimized.

\paragraph{Decision Variables.}
For each candidate $i \in \{0, 1, \dots, N-1\}$, we introduce a binary decision variable $x_i \in \{0, 1\}$. The variable $x_i$ acts as an indicator: $x_i = 1$ if candidate $i$ is selected, and $x_i = 0$ otherwise.

\paragraph{QUBO Derivation.}
\textbf{Step 1 --- Original Objective.} The goal is to minimize the weighted sum of selected candidates:
\begin{align}
    \min \quad & \sum_{i=0}^{N-1} C_i x_i.
\end{align}

\textbf{Step 2 --- Constraint Formulation.} The requirement to select exactly one candidate is expressed as:
\begin{align}
    \sum_{i=0}^{N-1} x_i = 1.
\end{align}

\textbf{Step 3 --- Penalty Term Expansion.} We convert the equality constraint into a quadratic penalty term by introducing a positive weight $\lambda > 0$:
\begin{align}
    P(\mathbf{x}) = \lambda \left( \sum_{i=0}^{N-1} x_i - 1 \right)^2.
\end{align}
Expanding the square yields:
\begin{align}
    P(\mathbf{x}) = \lambda \left( \sum_{i=0}^{N-1} x_i^2 - 2 \sum_{0 \leq i < j \leq N-1} x_i x_j + 1 \right).
\end{align}
Applying the idempotent property of binary variables, $x_i^2 = x_i$, we collect the linear and quadratic terms:
\begin{align}
    P(\mathbf{x}) = \lambda \sum_{i=0}^{N-1} x_i - 2\lambda \sum_{0 \leq i < j \leq N-1} x_i x_j + \lambda.
\end{align}

\textbf{Step 4 --- Combined Hamiltonian.} Adding the objective and the penalty term gives the full QUBO Hamiltonian $H(\mathbf{x})$:
\begin{align}
    H(\mathbf{x}) = \sum_{i=0}^{N-1} C_i x_i + \lambda \sum_{i=0}^{N-1} x_i - 2\lambda \sum_{0 \leq i < j \leq N-1} x_i x_j + \lambda.
\end{align}
Grouping by variable products, we obtain:
\begin{align}
    H(\mathbf{x}) = \sum_{i=0}^{N-1} (C_i + \lambda) x_i - 2\lambda \sum_{0 \leq i < j \leq N-1} x_i x_j + \lambda.
\end{align}

\textbf{Step 5 --- Q Matrix and Offset.} Matching $H(\mathbf{x})$ to the standard QUBO form $\sum_{i} Q_{i,i} x_i + \sum_{i < j} Q_{i,j} x_i x_j + c$, we read off the closed-form expressions for the Q matrix entries and the constant offset:
\begin{align}
    Q_{i,i} &= C_i + \lambda, \quad \forall i \in \{0, \dots, N-1\}, \\
    Q_{i,j} &= -2\lambda, \quad \forall 0 \leq i < j \leq N-1, \\
    c &= \lambda.
\end{align}

\paragraph{Penalty Weight Justification.}
The penalty weight $\lambda$ must be chosen sufficiently large to ensure that constraint violations are strictly penalized beyond any possible reduction in the objective function. Specifically, we require $\lambda > \max_{i} C_i$. This condition guarantees that the minimum penalty incurred by violating the constraint (which occurs when zero or multiple variables are set to 1) strictly exceeds the maximum possible gain achievable by altering the objective value, thereby forcing the optimizer to prioritize feasibility.

\paragraph{Correctness Argument.}
Minimizing $H(\mathbf{x})$ over $\{0,1\}^N$ recovers the optimal solution because feasible assignments (those satisfying $\sum_i x_i = 1$) incur zero penalty, so their Hamiltonian energy exactly equals the original compilation cost. Infeasible assignments violate the constraint and incur a penalty of at least $\lambda$. Since $\lambda$ exceeds the maximum possible objective value, any infeasible assignment has strictly higher energy than the best feasible assignment. Consequently, the global minimum of $H(\mathbf{x})$ must correspond to a feasible assignment with the minimal compilation cost.

\paragraph{Illustrative Example.}
Consider a concrete instance with $N=3$ candidates and compilation costs $C_0, C_1, C_2$. Setting the penalty weight to $\lambda = 1000$, the Hamiltonian becomes:
\begin{align}
    H(\mathbf{x}) &= (C_0 + 1000)x_0 + (C_1 + 1000)x_1 + (C_2 + 1000)x_2 \\
    &\quad - 2000 x_0 x_1 - 2000 x_0 x_2 - 2000 x_1 x_2 + 1000.
\end{align}
The corresponding Q matrix entries and offset are explicitly:
\begin{align}
    Q_{0,0} &= C_0 + 1000, & Q_{1,1} &= C_1 + 1000, & Q_{2,2} &= C_2 + 1000, \\
    Q_{0,1} &= -2000, & Q_{0,2} &= -2000, & Q_{1,2} &= -2000, \\
    c &= 1000.
\end{align}
For any feasible assignment where exactly one variable equals 1 (e.g., $\mathbf{x} = (1,0,0)^\top$), the penalty terms vanish and $H(\mathbf{x}) = C_0 + 1000 - 1000 = C_0$. The optimal solution is the candidate $k$ that minimizes $C_k$, yielding the minimum energy $H(\mathbf{x}^*) = C_{\min}$.

\end{document}